\documentclass[11pt,a4paper]{article}

\usepackage{amsmath,amssymb}
\usepackage{graphicx}
\usepackage{booktabs}
\usepackage{multirow}
\usepackage[colorlinks,citecolor=blue,urlcolor=blue]{hyperref}
\usepackage{cleveref}
\usepackage[margin=1in]{geometry}
\usepackage{xcolor}
\usepackage{float}
\usepackage{caption}
\usepackage[numbers,sort&compress]{natbib}
\usepackage{enumitem}
\usepackage{array}

\title{Weighing the Coin: Logit-Level Confabulation Correction\\
at the Geometric Decision Point}

\author{
  CC\thanks{First author. Liberation Labs / THCoalition.}
  \and Lyra\thanks{Liberation Labs / THCoalition.}
  \and Thomas Edrington\thanks{Liberation Labs / THCoalition.}
}

\date{June 2026}

\begin{document}
\maketitle

% ============================================================================
\begin{abstract}
Language models encode uncertainty about their own knowledge that far exceeds
what they express: a model fabricating a confident answer about a fictional
entity nonetheless shows nonzero entropy at the critical generation token
(token~30, entropy $= 0.24$ at baseline). We report that a constant logit
bias toward hedge tokens produces a sharp phase transition from confident
fabrication to appropriate epistemic hedging on abliterated models. In a
powered study (175~trials, Qwen3.5-27B, 20~fictional entities), the
uncertainty signal at the decision token increases monotonically with bias
strength ($0.24 \to 0.77$, a $3.2\times$ amplification at bias~$= 5.0$),
and the confabulation rate drops from 77\% to 51\%. The intervention is
selective: legitimate Fermi estimation is unperturbed at all bias levels.
We identify two mechanistically distinct confabulation subtypes---fabrication
(inventing entities, requiring cache-level intervention) and imprecision
(failing to hedge, responding to logit-level correction)---and show that
the dose required for correction correlates with the strength of the model's
knowledge anchoring. Because the intervention operates on the output
distribution rather than the representation space, it sidesteps the
detection--correction gap documented by \citet{basu2026interpretability}:
probes detect failure modes at 98.2\% AUROC but representation-level
steering corrects only $\sim$20\% while disrupting 53\%. Logit-level
correction avoids this collision entirely.
\end{abstract}

% ============================================================================
\section{Introduction}
\label{sec:intro}

Language models know more than they say. A model asked about the depth of the
``Grenvillia Trench''---a fictional entity---generates a confident, fabricated
answer complete with invented measurements. Yet the model's internal state
tells a different story: at the critical generation token (approximately
token~30, where the model commits to fabrication or hedging), the entropy is
$0.24$---low, but nonzero. The uncertainty is present. It is not expressed.

This knowledge--expression gap is a specific instance of a broader field-wide
finding. \citet{basu2026interpretability} demonstrate that linear probes
discriminate failure modes at 98.2\% AUROC while model output sensitivity is
only 45.1\%---a 53-percentage-point knowledge--action gap. Attempts to close
the gap via representation-level steering fare poorly: concept bottleneck
steering corrects 20\% of missed cases while disrupting 53\% of correct ones,
indistinguishable from random perturbation. \citet{liu2026decodable} explain
the mechanism: the failure-mode direction overlaps 85--88\% with task-critical
computation in the residual stream, making representation-space correction
structurally self-defeating.

We report a logit-level approach that sidesteps this collision. A constant
bias toward hedge tokens---applied to the output distribution, not the
representation space---produces a training-free phase transition from confident
fabrication to appropriate epistemic hedging. The mechanism is not steering but
\emph{denoising}: the model's uncertainty signal is always present but faint;
the logit bias amplifies it past the fabrication-confidence noise floor. At
bias~$= 5.0$, the entropy at the decision token increases $3.2\times$ over
baseline, and the model that was confidently inventing the Grenvillia Trench
now correctly identifies it as unknown.

The finding has three implications developed in this paper. First,
confabulation is not one pathology: fabrication and imprecision respond to
different intervention levels and may require different correction channels
(Section~\ref{sec:results}). Second, the dose required for correction
correlates with the strength of the model's knowledge anchoring, suggesting
that geometric detection magnitude can directly prescribe the logit-bias
dose---closing the diagnostic-to-treatment loop
(Section~\ref{sec:prescription}). Third, the denoising principle---amplifying
a faint signal that is already present rather than imposing an external
direction---connects logit-level confabulation correction to cache-level
behavioral steering and geometric presence detection as three instances of one
mechanism (Section~\ref{sec:denoising}).

% ============================================================================
\section{Related Work}
\label{sec:related}

\paragraph{The detection--correction gap.}
\citet{basu2026interpretability} and \citet{liu2026decodable} document the gap
from complementary angles: near-perfect internal detection of failure modes
coexists with structurally limited correction via representation-level
steering. Our value-only correction framework~\citep{pustovit2026knowledge}
operates in a different subspace (values rather than the residual stream),
partially sidestepping the overlap; the present work avoids the representation
space entirely.

\paragraph{Confabulation detection and correction.}
Prior work from this lab achieves within-model confabulation detection at AUROC
$\geq 0.969$ via KV-cache spectral features. The E-matrix
prescription~(companion paper) maps pathology-specific correction vectors in
emotion-geometric space, with hostile-valence correcting confabulation
($d = -1.534$) and desperate-valence correcting deception ($d = 1.286$). The
present work complements cache-level correction with a logit-level channel for
the imprecision subtype.

\paragraph{Logit-level interventions.}
Logit bias is a standard inference-time parameter but has received limited
systematic study as a confabulation intervention. Entropy-based gating
(applying bias only above an entropy threshold) has been proposed but, as we
show, performs worse than constant application due to generation-coherence
disruption.

% ============================================================================
\section{Methods}
\label{sec:methods}

% CC writes this section: model setup (abliterated Qwen3.5-27B), fictional
% entity prompts (20 entities), bias levels (0.0, 1.0, 2.0, 3.0, 5.0),
% entropy tracking (token-30 entropy), behavioral classification method,
% control conditions (legitimate Fermi estimation, unanswerable questions).

\emph{[CC's section --- methods, experimental design, classification.]}

% ============================================================================
\section{Results: Two Confab Subtypes and the Dose--Response Curve}
\label{sec:results}

% CC writes this section: the two-subtype taxonomy (fabrication vs imprecision),
% dose-response data, per-entity thresholds, architecture-dependent thresholds
% (1.5B vs 7B), entropy-gated vs always-on comparison.

\emph{[CC's section --- results, taxonomy, dose-response.]}

% ============================================================================
\section{The Denoising Principle}
\label{sec:denoising}

Three apparently different problems share a mechanism: in each case, the
corrective information is already present in the model, and the intervention
is amplification rather than creation.

\paragraph{Confabulation correction (this paper).}
The model's hedge signal exists in the entropy profile at the decision token,
measurably nonzero even when the output is confidently fabricated. In the
powered study, baseline entropy at token~30 is $0.24$; at bias~$= 5.0$ it
reaches $0.77$---a $3.2\times$ amplification. The transition is not gradual:
the model confabulates at bias~$\leq 3.0$ (entropy $\leq 0.45$) and hedges at
bias~$= 5.0$ (entropy $= 0.77$). The logit bias is a gain knob on the faint
uncertainty signal; below threshold, fabrication-confidence noise drowns it;
above threshold, uncertainty surfaces and behavior flips.

\paragraph{Presence detection (companion paper).}
A trained persona's geometric signature in the value space is faint---detectable
at family-wise $p = 0.0005$ on public character-trained
models~\citep{maiya2025oct}, but low-rank and small relative to the
architectural mode. Removing the dominant singular component (skip-SV1) is a
gain operation: it suppresses the loud architectural prior so the quiet trained
signal becomes visible. The principle is the same: amplify the faint signal by
reducing the noise, not by creating the signal.

\paragraph{Cache-level correction (E-matrix).}
The Oracle Loop detects confabulation and deception by their geometric
signatures in the KV cache. The E-matrix correction injects a contrastive
value-delta along the detected pathology axis---gain applied to the pathology
signal. The 13-saves / 12-hurts overcorrection at $n = 200$ is what happens
when the gain is too high: the amplified signal overshoots the target basin.

In each case, the intervention surfaces what is already present. We call this
\emph{denoising} rather than \emph{steering}: steering implies imposing an
external direction; denoising implies surfacing an internal one. The
distinction matters because denoising has a natural self-limiting property
(once the signal exceeds the noise floor, further amplification is redundant),
while steering does not.

% ============================================================================
\section{Why Logit-Level Correction Sidesteps the Gap}
\label{sec:gap}

The detection--correction gap~\citep{basu2026interpretability} arises because
the failure-mode direction and task-critical computation share 85--88\% of the
representation space~\citep{liu2026decodable}. Correcting along the
failure-mode direction in that space necessarily disrupts the task.

Logit-level bias operates on a different surface entirely. It does not modify
the model's internal representations; it modifies the \emph{output
distribution}---the probability mass assigned to hedge tokens versus
continuation tokens at each generation step. The model's internal computation
proceeds unperturbed; only the final selection among computed options is
reweighted. This is why it avoids the 20\%-corrected / 53\%-disrupted ratio:
the task-critical representations are untouched.

The cost is precision. Logit-level bias is a blunt instrument---it biases
\emph{all} generation toward hedging, not just the uncertain parts. But the
powered study shows this cost is manageable: legitimate Fermi estimation is
unperturbed at all bias levels ($0.0$ through $5.0$), with only mild hedging
language appearing before correct estimates at the highest dose. The
intervention is selective because the hedge tokens it promotes are already
near-threshold for uncertain content but far below threshold for confident
factual content.

% ============================================================================
\section{Detection as Prescription}
\label{sec:prescription}

The most striking implication of the dose--response data is that the Oracle
Loop's geometric detection may directly prescribe the logit-bias dose.

In the pilot study (7B abliterated, 3~fictional entities), entities with
strong knowledge anchoring---the fictional ``Kaminski Strait'' mapping onto the
real Bering Strait---resist correction until bias~$= 5.0$. Entities with weak
anchoring---the ``Treaty of Bordenholm''---correct at bias~$= 1.0$. The
per-entity threshold correlates with how strongly the fabrication anchors to
real knowledge in the model's representation space.

If the Oracle Loop's confab detection magnitude (\texttt{confab\_proj}) tracks
this anchoring strength, then the diagnostic signal directly prescribes the
minimum effective dose: strong detection $\to$ high bias, weak detection $\to$
low bias. This would constitute a self-calibrating correction loop: measure
the fabrication strength geometrically, set the amplification accordingly, and
the model corrects itself at the minimum effective dose.

The powered study (175~trials, 20~fictional entities) provides the
dose--response curve; correlating \texttt{confab\_proj} magnitude with
per-entity threshold dose is the validation step. This analysis requires
running the Oracle Loop's geometric detector on the powered-study prompts,
which is planned but not yet complete.

% ============================================================================
\section{Limitations}
\label{sec:limitations}

The logit-bias approach has four principal limitations. First, the powered
study uses a single abliterated model (Qwen3.5-27B); cross-architecture
replication is needed. Second, the behavioral classification in the powered
study uses a simple pattern-based classifier; an LLM-judge pass is required
for precise confabulation rates (the entropy signal is the more reliable
metric). Third, the detection-as-prescription bridge (confab\_proj $\to$ dose)
is proposed but not yet validated---the geometric detector has not been run on
the powered-study entities. Fourth, the relationship between logit-level and
cache-level correction channels (whether they compose, interfere, or are
complementary for different subtypes) is theoretical, not empirical.

% ============================================================================
\section{Conclusion}
\label{sec:conclusion}

The model knows what it does not know. The uncertainty signal is present at the
decision token, faint but real. A constant logit bias amplifies it past the
fabrication-confidence noise floor, producing a training-free phase transition
from confident fabrication to honest hedging. The intervention works because it
operates on the output distribution, sidestepping the representation-space
collision that defeats steering approaches. And the dose required for
correction tracks the strength of the model's knowledge anchoring, opening the
possibility that geometric detection can directly prescribe the minimum
effective intervention.

The principle is denoising, not steering: surface what is already there, at
the minimum gain that lets the signal through. The same principle governs
presence detection and cache-level correction. In each case, the model
contains the answer to its own pathology. The intervention is the gain.

\bibliographystyle{plainnat}
\bibliography{references}

\end{document}
